\section{Syntax}\label{sec:syntax}

\subsection{Well-typed formulas}\label{subsec:wtform}

\subsubsection{Formal expressions}\label{ssubsec:formex}
Let $0\leq k\leq n$, $C$ be an $\pair nk$-category, and  $\clx CS$ be a
cellular extension of $C$ by a set $S$, with $\sce n^S,\tge n^S:S\to C_n$. For each $a\in S$, and $x$,
$y\in C_n$, the notation $a:x\to y$ is short for 
$x=\sce n^S(a)$ and $y=\tge n^S(a)$. To $\clx CS$ we first associate
the formal language $\expr S$ whose formulas are given by the grammar
\[
  e ::= \fcst{a} \mid \finv{a} \mid \fid{x} \mid (e\fcomp i e)
\]
where $a\in S$, $x\in C_n$ and $0\leq i\leq n$. Note that this
language has the unique parsing property. This allows reasoning by
structural induction on the complexity of formulas in $\expr S$. Thus,
in order to completely determine a map $f$ from $\expr S$ to an arbitrary
set $X$, it suffices to define $f(e)$ for $e$ of the form $\fcst a$,
$\finv a$ or $\fid x$ and to provide the rules to evaluate $f((e_1\fcomp
i e_2))$ from the values $f(e_1)$ and $f(e_2)$ for each $i\leq n$.
For instance, we may define a {\em length} map $\lgh:\expr S\to \N$
by setting $\lgh \fcst a=\lgh\finv a= \lgh \fid x=1$ and $\lgh
(e_1\fcomp i e_2)=\lgh e_1+\lgh e_2$. 



\subsubsection{Typing}\label{ssubsec:types}
We now define the
subset  $\wtx S\subset\expr S$ of {\em well-typed formulas}.
Precisely, we define the statement
\[\type wxy\]
meaning that {\em $w\in\wtx S$ is of type $\pair xy$} with $x$, $y$ parallel
cells in $C_n$, by the following structural induction:
\begin{itemize}
\item for each $a:x\to y\in S$ ,
  $\type{\fcst{a}}{x}{y}$
  and $\type{\finv{a}}{y}{x}$;
\item for  each $x\in C_n$, $\type{\fid{x}}{x}{x}$;
 \item if $\type wxy$ and $\type{w'}yz$, then $w''=(w\fcomp n w')$ is
   well-typed and $\type{w''}{x}{z}$;
 \item if $0\leq i<n$, $\type wxy$, $\type{w'}{x'}{y'}$, $\sce
   i{x'}=\tge i{x}$, then $w''=(w\fcomp i w')$ is well-typed and
   $\type{w''}{x\comp i x'}{y\comp i y'}$. 
 \end{itemize}
 One readily checks that whenever $\type wxy$ is derivable from the
 preceding rules, then $x$, $y$ are well defined parallel cells in
 $C_n$, so that the maps $\sce n^S$, $\tge n^S$ extend to maps
 \[\scex n^S, \tgex n^S :\wtx S\to C_n\]
 such that, if $\type wxy$,  then $\scex n^S w=x$ and $\tgex n^S
 w=y$. The map $\finc S:S\to \wtx S$  taking
 $a\in S$ to $\fcst a\in\wtx S$ factorizes $\sce n^S$ and $\tge n^S$
 as $\sce n^S=\scex n^S \finc S$ and $\tge n^S=\tgex n^S\finc S$
 respectively. By composition with
 $\sce i^C,\tge i^C:C_n\to C_i$ for any $i<n$, we define likewise
 $\scex i^S,\tgex i^S:\wtx S\to C_i$. Note that the map $\finc S$ is injective. 
 
 \begin{paragr}\label{paragr:transl}
  Let $C$, $D$ be $\pair nk$-categories, $S$, $T$ be two sets,  $\clx
  CS$, $\clx DT$ be cellular extensions and
  \[\phi=\pair fh: \clx CS\to\clx DT \]
  be a morphism in $\npCatp nk$. A map $g:\wtx S\to\wtx T$ is called a {\em
   $\phi$-compatible translation} if
  \begin{itemize}
  \item for each $w\in\wtx S$, $\scex n^Tg(w)=f_n(\scex n^S w)$ and
    $\tgex n^T g(w)=f_n(\tgex n^S w)$;
    \item for each $a\in S$, $g(\fcst a)=\fcst{h(a)}$ and $g(\finv
      a)=\finv{h(a)}$;
    \item for each $x\in C_n$,  $g(\fid x)=\fid{f_n(x)}$;
    \item  for each $w_1,w_2\in\wtx S$ and $i\leq n$ such that
      $w=(w_1\fcomp i w_2)$ is well-typed, $g(w)=(g(w_1)\fcomp i g(w_2))$.  
    \end{itemize}
\end{paragr}
\begin{lemma}\label{lemma:transl}
  Let $\phi=\pair fh:\clx CS\to \clx DT$ be a morphism of cellular
  extensions. There is a unique $\phi$-compatible translation $\wtx\phi:\wtx S\to \wtx T$.
\end{lemma}
\begin{proof}
  By structural induction on $w\in\wtx S$, there is a unique map
  $g:\wtx S\to \expr T$  such that $g(\fcst a)=\fcst{h(a)}$, $g(\finv
      a)=\finv{h(a)}$, $g(\fid x)=\fid{f_n(x)}$ and $g((w_1\fcomp i
      w_2))=(g(w_1)\fcomp i g(w_2))$. For this map $g$, we easily
      show, by structural induction on $w$, that
      whenever $\type wxy$, then $\type{g(w)}{f_nx}{f_n y}$. Therefore
      $g$ factors through the subset $\wtx T\subset \expr T$ of
      well-typed formulas and yields the required $\phi$-compatible
      translation $\wtx\phi:\wtx S\to \wtx T$. 
    \end{proof}
    \begin{remark}\label{rmk:functoriality_of_wtx}
      By \cref{lemma:transl}, the correspondence
      \[
        \clx CS\mapsto \wtx S, (\phi:\clx CS\to \clx
      DT)\mapsto \wtx\phi
        \]
     defines a functor 
      $\WW:\npCatp nk\to\Set$.
    \end{remark}

 \begin{paragr}
   Let $\clx CS$ be a cellular extension of the $\pair nk$-category
   $C$ by a set $S$, $D$ be an $\pair{n{+}1}k$-category and $f:C\to \trk n
   D$ be a morphism. A map $g:\wtx S\to D_{n+1}$ is called {\em $f$-compatible} if:
   \begin{itemize}
   \item for each $w\in\wtx S$, $\sce n^Dg(w)=f_n(\scex n^S w)$ and
     $\tge n^D g(w)=f_n(\tgex n^S w)$;
    \item for each $a\in S$, $g(\finv a)$ is the inverse of $g(\fcst
      a)$;
    \item for each $x\in C_n$, $g(\fid x)=\unit{n+1}(f_n(x))$;
    \item for each $w_1,w_2\in\wtx S$ and $i\leq n$ such that
      $w=(w_1\fcomp i w_2)$ is well-typed, $g(w)=g(w_1)\comp i g(w_2)$.  
    \end{itemize}
    Given $C$, $S$, $D$ and $f$ as above, we obtain the following result.
  \end{paragr}
  \begin{lemma}\label{lemma:mapext}
    Let $h:S\to D_{n+1}$ be a map such that $\sce n^Dh=f_n \sce
    n^S$ and $\tge n^Dh=f_{n} \tge n^S$. There is a unique
    $f$-compatible map
    $\extend h:\wtx S\to D_{n+1}$
    such that $\extend h\finc S=h$.
  \end{lemma}
  \begin{proof}
    Uniqueness immediately follows from $f$-compatibility. As for the
    existence, we define, for each $w:x\to y\in\wtx S$,
  $\extend h(w):f_n(x)\to f_n(y)\in D_{n+1}$, by structural induction
  on $w$:
  \begin{itemize}
  \item for $w=\fcst a$, where $a:x\to y\in S$, $\extend
    h(w)=h(a):f_n(x)\to f_n(y)$;
  \item for $w=\finv a$, where $a:x\to y\in S$, $\extend
    h(w)=\inv{h(a)}:f_n(y)\to f_n(x)$;
  \item for $w=\fid x$, where $x\in C_n$, $\extend
    h(w)=\unit{n+1}(f_n(x))$;
  \item if $w_1:x\to y$ and $w_2:y\to z$ are such that $\extend
    h(w_1):f_n(x)\to f_n(y)$ and $\extend h(w_2):f_n(y)\to f_n(z)$,
    then $w=(w_1\fcomp n w_2):x\to z$ and $d=\extend h(w_1)\comp
    n\extend h(w_2)$ is a well defined cell of $D_{n+1}$ such that
    $d:f_n(x)\to f_n(z)$ and we may set $\extend h(w)=d$;
  \item if $i<n$, $w_1:x\to y$, $w_2:z\to t$, $\sce i^Cz=\tge i^Cx $
    are such that $\extend h(w_1):f_n(x)\to f_n(y)$ and
    $\extend h(w_2):f_n(z)\to f_n(t)$, then
    $w:x\comp i z\to y\comp i t\in \wtx S$ and $\sce i^D\extend
    h(w_2)=\sce i^Df_n(z)=f_i(\sce i^Cz)=f_i(\tge i^Cx)=\tge
    i^Df_n(x)=\tge i^D\extend h(w_1)$, so that $d=\extend h(w_1)\comp
    i\extend h(w_2)$ is a well defined cell in $D_{n+1}$ with
    $d:f_n(x\comp i z)\to f_n(y\comp i t)$ and we may set $\extend h(w)=d$.
  \end{itemize}
  By construction, the map $\extend h$ so defined is $f$-compatible
  and $\extend h\finc S =h$.
  \end{proof}
  
  
\subsection{Formal description of $\frf {n,k}$}\label{subsec:congruence}

  
 \subsubsection{Equivalence among formal
   expressions}\label{ssubsec:eqform}
 We now define the binary relation of ``elementary move'' among well
 typed expressions, denoted by $\frel$, as follows:
\begin{enumerate}[label=\textbf{(e\arabic*)}]
\item\label{ei} for any pair $x$, $x'$ of $i$-composable cells in $C_n$,
  $(\fid{x}\fcomp i\fid{x'})\frel\fid{x\comp i x'} $;
\item\label{eii} for each $a:x\to y$ in $S$, $(\fid{x}\fcomp{n}\fcst{a})\frel
  \fcst{a}$ and $(\fcst a\fcomp{n} \fid{y})\frel \fcst a$;
 \item\label{eiii} for each $a:x\to y$ in $S$, $(\fid{y}\fcomp{n}\finv{a})\frel
   \finv{a}$ and $(\finv a\fcomp{n} \fid{x})\frel \finv a$;
\item\label{eiv} for each $a:x\to y$ in $S$, $0\leq i< n$, $u=\unit n(\sce i
  x)$ and $v=\unit n(\tge i x)$, $(\fid{u}\fcomp{i}\fcst{a})\frel
  \fcst{a}$ and $(\fcst a\fcomp{i} \fid{v})\frel \fcst a$;
\item\label{ev} for each $a:x\to y$ in $S$, $0\leq i< n$, $u=\unit n(\sce i
  x)$ and $v=\unit n(\tge i x)$, $(\fid{u}\fcomp{i}\finv{a})\frel
  \finv{a}$ and $(\finv a\fcomp{i} \fid{v})\frel \finv a$;  
 \item\label{evi} for all expressions $w_1,w_2,w_3$ in $\wtx S$ such that $w_1,w_2$ and
   $w_2,w_3$ are $i$-composable, $((w_1\fcomp i w_2)\fcomp i w_3)\frel
   (w_1\fcomp i(w_2\fcomp i w_3))$;
 \item\label{evii} for all expressions $w_1,w_2,w_3,w_4$ in $\wtx S$ and $i<j$
   such that $w_1,w_2$ and $w_3,w_4$ are
   $j$-composable and $w_1,w_3$ are $i$-composable,
   \[
     ((w_1\fcomp i w_3)\fcomp j (w_2\fcomp i w_4))\frel ((w_1\fcomp j
     w_2)\fcomp i (w_3\fcomp j w_4));
   \]
  \item\label{eviii} for each $a:x\to y$ in $S$, $(\fcst a\fcomp n\finv a)\frel
    \fid{x}$ and $(\finv a\fcomp n\fcst a)\frel \fid y$. 
  \end{enumerate}
  Let now $\freq$ denote the congruence generated by
  $\frel$ on $\wtx S$, that is, the smallest equivalence relation
  compatible with compositions containing $\frel$. Precisely, if
  $w,w'\in\wtx S$, the relation  ``$w$ rewrites in one step
  to $w'$'',
  denoted by $w\cfrel w'$, is defined inductively as follows: $w\cfrel
  w'$ if either (i)~$w\frel w'$ or (ii)~$w=(x\fcomp i u)$ and $w'=(x\fcomp i u')$
  with $u\cfrel u'$ or (iii)~$w=(u\fcomp i x)$,
  and $w'=(u'\fcomp i x)$ with $u\cfrel u'$. The relation $\freq$ is
  then the smallest equivalence relation containing $\cfrel$. In other
  words, $w\freq w'$ if and only if there is a finite sequence
  $w_0,...,w_l$ where $l\geq 0$ of formulas in $\wtx S$ such that
  $w=w_0$, $w'=w_l$ and for each $i<l$, either $w_i\cfrel w_{i+1}$ or
  $w_{i+1}\cfrel w_i$.

  Let us denote by $\wtxeq S$
  the set $\wtx S/\!\!\freq$ of equivalence classes of $\wtx S$ under
  $\freq$ and by
  \[\cans S: \wtx S\to \wtxeq S\]
  the canonical surjection taking an expression $w$ to its
  equivalence class $\cans S(w)=\fcls w$. As the type of an
  expression is invariant under $\frel$, therefore also under $\freq$, 
 the source and target maps
  \[\scex n^S,\tgex n^S:\wtx S\to C_n\]
  factor through $\cans S$, and define source and target maps on $\wtxeq S$:
  \[\scecl n^S,\tgecl n^S:\wtxeq S\to C_n.\]
  More generally, for any $i\leq n$, we define the $i$-source and
  $i$-target maps
  \[\scecl i,\tgecl i:\wtxeq S\to C_i\]
  by composing $\sce n^S,\tge n^S:\wtxeq S\to C_n$ with the maps $\sce i^C,\tge
  i^C:C_n\to C_i$ already given. Now, for any $i$-composable formal expressions
  $w,w'\in \wtx S$ and $w''=(w\fcomp i w')$, the class $\cans S(w'')=\fcls{w''}$ only
  depends on $\fcls w$ and $\fcls{w'}$. Hence we may define 
 $\comp i$-compositions in $\wtxeq S$ by
 \[\fcls{w}\comp i\fcls{w'}=\fcls{(w\fcomp i w')}.\]
 Also, to any $x\in C_n$ we associate its $n{+}1$-unit in $\wtxeq
  S$ by
  \[\unit{n+1}(x)=\fcls{\fid x}.\]
  \begin{lemma}\label{lemma:inv}
    For each $\type wxy\in \wtx S$, there is a $\type{w'}yx\in \wtx S$
    such that $(w\fcomp n w')\freq\fid x$ and $(w'\fcomp n w)\freq
    \fid y$.
  \end{lemma}
  \begin{proof}
    We reason by induction on the complexity of $w$:
    \begin{itemize}
    \item if $w=\fid u$ for some $u\in C_n$, then $x=y=u$ and we set
      $w'=w$;
    \item if $w=\fcst a$ for some $a\in S$, we set $w'=\finv a$;
    \item if $w=\finv a$ for some $a\in S$, we set $w'=\fcst a$;
     \item if $w=(w_1\fcomp n w_2)$, then $\type{w_1}xz$ and
       $\type{w_2}zy$ and by induction we get $\type{w'_1}zx$ and
       $\type{w'_2}yz$ such that $(w_1\fcomp n w'_1)\freq \fid x$,
       $(w'_1\fcomp n w_1)\freq \fid z$, $(w_2\fcomp n w'_2)\freq \fid
       z$
       and $(w'_2\fcomp n w_2)\freq \fid y$. Thus $w'=(w'_2\fcomp n
       w'_1)$ is well typed and satisfies $(w\fcomp n w')\freq \fid x$
       and $(w'\fcomp n w)\freq \fid y$;
       \item if $w=(w_1\fcomp i w_2)$ for some $i<n$, then
         $\type{w_1}{x_1}{y_1}$
         and $\type{w_2}{x_2}{y_2}$ with $x=x_1\comp i x_2$ and
         $y=y_1\comp i y_2$. By induction, we get
         $\type{w'_1}{y_1}{x_1}$ and $\type{w'_2}{y_2}{x_2}$ such that
         $(w_l\fcomp n w'_l)\freq \fid{x_l}$ and $(w'_l\fcomp n
         w_l)\freq \fid{y_l}$ for $l=1,2$. Thus $w'=(w'_1\fcomp i
         w'_2)$ is well typed and satisfies $(w\fcomp n w')\freq \fid x$
       and $(w'\fcomp n w)\freq \fid y$.
    \end{itemize}
  \end{proof}
  Let us define the $n{+}1$-globular set $\extend C$ by
  $\trk n\extend C=C$ and $\sce n,\tge n:\extend C_{n+1}=\wtxeq S\to C_n$.
  The axioms for $\frel$ ensure that $\extend C$, endowed with
  the $\comp i$-compositions and units just defined is an
  $(n{+}1)$-category, and by~\cref{lemma:inv}, all its $(n{+}1)$-cells
  are invertible, so that $\extend C$ is
  indeed an
  $\pair{n{+}1}{k}$-category.
 

  We claim that $\extend C$ so defined is
  in fact, up to isomorphism, the expected $\pair{n{+}1}{k}$-category
  $\frf{n,k}\clx CS$ freely generated by the cellular extension $\clx CS$.
  This amounts to prove that $\extend C$ satisfies the following universal property.
  
\begin{proposition}\label{prop:univprop}
   Let $D$ be an  $\pair{n{+}1}{k}$-category, $f:C\to\trk n D$ be a
   morphism of $\pair nk$-categories
   and  $h:S\to D_{n+1}$ be a map such that, for any $a\in S$, $\sce n^Dh(a)=\sce n^S a$
    and $\tge n^D h(a)=\tge n^S a$.  There is a unique
    morphism $\phi:\extend C\to D$ in $\npCat{n+1}{k}$ such that $\trk
    n\phi=f$ and for each $a\in S$, $\phi_{n+1}(\fcls{\fcst a})=h(a)$.
  \end{proposition}
  \begin{proof}
We first prove uniqueness. As $\phi$ is supposed to coincide with $f$ on
dimensions $\leq n$, we only need to prove uniqueness for its
$(n{+}1)$-component
\[\phi_{n+1}:\wtxeq S\to D_{n+1}.\]
 Thus, suppose
that $\phi$ satisfies the conditions of the proposition and let
\[\psi=\phi_{n+1} \cans{S}:\wtx S\to D_{n+1}.\]
   Because $\phi$ is a morphism, the  map $\psi$ is
   $f$-compatible, and by hypothesis $\psi(\fcst
   a)=\phi_{n+1}\cans S(\fcst a)=h(a)$. Therefore, by the uniqueness
   part of
   Lemma~\ref{lemma:mapext}, the map $\psi=\extend h$ is uniquely determined.
As $\cans S$  is surjective,
  $\phi_{n+1}$ is also uniquely determined and so is $\phi$.

  As for the existence, by Lemma~\ref{lemma:mapext}, there is an
  $f$-compatible map
  \[\psi:\wtx S\to D_{n+1}\]
  such that $\psi(\fcst a)=h(a)$ for each $a\in S$. 
  A straightforward case analysis shows that whenever $w\frel w'$,
  then $\psi(w)=\psi(w')$. Therefore, if $w\freq w'$, then
  $\psi(w)=\psi(w')$, whence $\psi$ factors through $\cans S$, that
  is, there is a map $\phi_{n+1}:\wtxeq S\to D_{n+1}$ such that
  $\psi=\phi_{n+1}\cans S$. This map preserves source and target,
  identities and $\comp i$-compositions and, by definition, $\phi_{n+1}(\fcls
  {\fcst{a}})=h(a)$, so that it extends $f$ to the required morphism $\phi$. 
\end{proof}
\begin{remark}\label{rmk:functoriality_of_wtxeq}
  By~\cref{prop:univprop}, the correspondence
  \[\clx CS\mapsto \wtxeq S\]
  is the object part of a functor $\clWW:\npCatp
  nk\to\Set$. Precisely, this functor $\clWW$ is the composite
  \[
    \xymatrix{\npCatp nk\ar[r]^{\frf{n,k}} & \npCat{n+1}k \ar[r]^(.6){(-)_{n+1}}& \Set}
  \]
  where $(-)_{n+1}$ takes an $\pair{n{+}1}k$-category to the set of
  its $(n{+}1)$-cells.

  Let now $\phi=\pair fh:\clx CS\to \clx DT$
  be a morphism of cellular extensions, and define $h':S\to\wtxeq T$
  by $h'(a)=\fcls{\fcst{h(a)}}$. By~\cref{prop:univprop},
  $\extend\phi=\frf{n,k}\phi:\frf{n,k}\clx CS\to\frf{n,k}\clx DT$ is
  the unique morphism in $\npCat{n+1}k$ such that $\trk
  n\extend\phi=f$ and $\extend{\phi}_{n+1}(\fcls{\fcst a})=h'(a)$ for
  each $a\in S$. Now
  $\extend{\phi}_{n+1}$ is precisely $\wtxeq{\phi}$, so that, for each $a\in
  S$, $\wtxeq{\phi}(\fcls{\fcst a})=\fcls{\fcst{h(a)}}$.
\end{remark}

\begin{paragr}\label{paragr:addproperties}
  Let $C$ be an $\pair{n{+}1}{k}$-category and $S$ be a set. Any map
  $p:S\to C_{n+1}$ immediately defines a cellular extension $\clx{\trk n C}{S}$ of
  $\trk n C$ by $S$, with $\sce n^S=\sce n^C\circ p$ and $\tge
  n^S=\tge n^C\circ p$. Thus
\[\pair{\Unit{\trk n C}}{p}:\clx{\trk n C}{S}\to\fgf{n,k}C\]
is a morphism of cellular extensions. Let $\extend C=\frf{n,k}\clx{\trk n C}{S}$ and $\phi:\extend C\to C$
be the transpose of $\pair{\Unit{\trk n C}}{p}$ under the adjunction. We know
by~\cref{prop:univprop} that $\extend{C}_{n+1}=\wtxeq S$. Let $\fcls
p=\phi_{n+1}:\wtxeq S\to C_{n+1}$ and $\extend p=\fcls p\cans S:\wtx
S\to C_{n+1}$. By~\cref{lemma:mapext}, the map $\extend p$ is the
unique $\Unit{\trk nC}$-compatible map such that $\extend p\finc
S=p$. Therefore, each $p:S\to C_{n+1}$ yields a factorization of the
form
\begin{equation}
  \label{eq:factor}
  \xymatrix{C_{n+1} & \wtxeq S \ar[l]^{\fcls{p}}& \wtx S\ar[l]^(.4){\cans{S}} \ar@/_1em/[ll]_(.3){\extend{p}}& S\ar[l]^(.3){\finc{S}}\ar@/_2em/[lll]_p}
\end{equation}
\end{paragr}


\begin{lemma}\label{lemma:main_diagram}
    Let $C$, $D$ be $\pair{n{+}1}{k}$-categories and $f:C\to D$ be a
    morphism. Let $S$, $T$ be two sets endowed with maps $p:S\to
    C_{n+1}$ and $q:T\to D_{n+1}$. Let $h:S\to T$ be a map such that
    $f_{n+1}p=qh$. Then
    \[ \phi=\pair{\trk n f}{h}:\clx{\trk n
      C}S\to\clx{\trk n D}{T}\]
is a morphism of cellular extensions
    and the following diagram commutes:
    \begin{equation}
      \label{eq:main_diagram}
      \begin{tikzcd}
       C_{n+1} \ar[d,"f_{n+1}"']& \wtxeq S\ar[l,"\fcls{p}"]\ar[d,"\wtxeq\phi"'] & \wtx S
      \ar[l,"\cans{S}"]\ar[d,"\wtx\phi"] & S\ar[l,"\finc{S}"]\ar[d,"h"]
      \ar[lll,bend right=15,"p"']\\
      D_{n+1} & \wtxeq T\ar[l,"\fcls{q}"'] & \wtx T
      \ar[l,"\cans{T}"'] & T.\ar[l,"\finc{T}"'] \ar[lll,bend left=15,"q"]
      \end{tikzcd}
     % \xymatrix{
     % C_{n+1} \ar[d]_{f_{n+1}}& \wtxeq S\ar[l]^{\fcls{p}}\ar[d]_{\wtxeq\phi} & \wtx S
     %  \ar[l]^(.4){\cans{S}}\ar[d]^{\wtx\phi} & S\ar[l]^(.3){\finc{S}}
     %  \ar@/_1em/[lll]_p\ar[d]^h\\
     %  D_{n+1} & \wtxeq T\ar[l]_{\fcls{q}} & \wtx T
     %  \ar[l]_(.4){\cans{T}} & T\ar[l]_(.3){\finc{T}} \ar@/^1em/[lll]^q}.
     \end{equation}
   \end{lemma}
   \begin{proof}
     The outer square commutes by hypothesis, and the top and bottom
     factorizations are instances
     of~(\ref{eq:factor}). By~\cref{lemma:transl}, $\wtx\phi$ is a
     $\phi$-compatible translation and in particular, for each $a\in
     S$, $\wtx\phi(\fcst a)=\fcst{h(a)}$ so that the right square
     commutes.

     Let $h'=\cans T\finc Th:S\to\wtxeq T$. Both maps
     $\wtxeq\phi\cans S$ and $\cans T\wtx\phi$ are
     $f$-compatible. Moreover, for each $a\in S$,
\[\wtxeq\phi\cans S\finc S(a)=\wtxeq\phi(\fcls{\fcst a})=\fcls{\fcst{h(a)}}=h'(a)\]
by~\cref{rmk:functoriality_of_wtxeq} and
\[
\cans
     T\wtx\phi\finc S(a)=\cans T\finc T h(a)=h'(a).
  \]
Thus, by the
     uniqueness part of~\cref{lemma:mapext}, $\wtxeq\phi\cans S=\cans
     T\wtx\phi$, so that the middle square commutes.

     Let finally
     $h''=qh$. Both maps $f_{n+1}\fcls p\cans S$ and $\fcls
     q\wtxeq\phi\cans S$ from $\wtx S$ to $D_{n+1}$ are $f$-compatible. Also
     \[
       f_{n+1}\fcls p\cans S\finc S=f_{n+1}p=qh=h''
     \]
      and
      \[
      \fcls q\wtxeq\phi\cans S\finc S=\fcls q\cans T\wtx\phi\finc
      S=\fcls q\cans T\finc T h=qh=h''.
      \]
      By the uniqueness part
     of~\cref{lemma:mapext},
     $f_{n+1}\fcls p\cans S=\fcls q\wtxeq\phi\cans S$ and as $\cans S$
     is surjective, $f_{n+1}\fcls p=\fcls q\wtxeq\phi$. Therefore the
     left square commutes and we are done.
     
   \end{proof}

%%% Local Variables:
%%% mode: LaTeX
%%% TeX-master: "main"
%%% End:
